\documentclass{article}

% Language setting
% Replace `english' with e.g. `spanish' to change the document language
\usepackage[english]{babel}

% Set page size and margins
% Replace `letterpaper' with `a4paper' for UK/EU standard size
\usepackage[letterpaper,top=2cm,bottom=2cm,left=3cm,right=3cm,marginparwidth=1.75cm]{geometry}

% Useful packages
\usepackage{amsmath}
\usepackage{graphicx}
\usepackage[colorlinks=true, allcolors=blue]{hyperref}

\title{
    Rapport mi-parcours
    \\ \medskip \medskip \small Projet Industriel MAIN3
    }
    
\author{Léa  \\ Alma  \\ Khireddine  \\ Lamia}
\date{30 Mars 2026}

\begin{document}
\maketitle

\begin{abstract}
Un résumé ? 
\emph{
\\ Dans le résumé, on pourrait donner
    \begin{itemize}
        \item indiquer le sujet et une explication / notre interprétation ?
        \item donner une définition de LLM, RAG, etc...
    \end{itemize} 
}
\end{abstract}

\section{Introduction}

On a obtenu plus d'infos à l'aide du livre \cite{livre_RAG}. 
\\
Qu'est-ce que le embedding ? \cite{embedding1}
% je mets ça ici juste pour le faire apparaître dans la biblio :) 


\section{Contexte du travail}


\section{Cahier des charges}

\section{Etat de l'art}

\section{Description des réalisations préliminaires}

\section{Etat de l'avancement}

% ===== ICI LA BIBLIOGRAPHIE =====
\newpage

\begin{thebibliography}{9}
\addcontentsline{toc}{section}{Références}

    \bibitem{livre_RAG}
        Essential GraphRAG - KNOWLEDGE GRAPH-ENHANCED RAG : 
        \href{Essential-GraphRAG.pdf}{Essential-GraphRAG.pdf}

    \bibitem{embedding1}
        Embedding Models Explained: A Guide to NLP’s Core Technology
        \href{https://medium.com/@nay1228/embedding-models-a-comprehensive-guide-for-beginners-to-experts-0cfc11d449f1}{https://medium.com/@nay1228/embedding-models-a-comprehensive-guide-for-beginners-to-experts-0cfc11d449f1}



\end{thebibliography}


\end{document}